\begin{apendicesenv}

% Imprime uma página indicando o início dos apêndices
\partapendices

% ===========================================================================
\chapter{Especificação Detalhada dos Módulos e Gatilhos GTT}
\label{apend_modulos_gatilhos}
% ===========================================================================

O Instituto para Melhoria da Saúde (IHI) definiu 53 gatilhos clínicos
distribuídos em sete módulos assistenciais para a metodologia
\textit{Global Trigger Tool} (GTT) \cite{griffin2009ihi}. O
\autoref{qua_modulos_completos} apresenta a especificação completa
dos módulos e respectivos gatilhos implementados no catálogo do SIGEA-GTT.

\begin{quadro}[htb]
\centering
\small
\caption{Especificação dos Módulos Assistenciais e Gatilhos GTT (IHI).}
\label{qua_modulos_completos}
\begin{tabularx}{\textwidth}{l l X}
\toprule
\textbf{Módulo} & \textbf{Código} & \textbf{Descrição do Gatilho} \\
\midrule
\multirow{10}{*}{\textbf{Cuidados Medicamentosos (CM)}}
  & CM-1 & Diphenhydramine (Benadryl) ou outro antihistamínico \\
  & CM-2 & Vitamina K administrada \\
  & CM-3 & Flumazenil (Romazicon) administrado \\
  & CM-4 & Naloxona (Narcan) administrada \\
  & CM-5 & Antieméticos administrados \\
  & CM-6 & Superdosagem de anticoagulante oral \\
  & CM-7 & Hipoglicemia registrada \\
  & CM-8 & Subida de creatinina $\geq 2\times$ basal \\
  & CM-9 & INR $> 6$ registrado \\
  & CM-10 & Trombocitopenia pós-heparina \\
\midrule
\multirow{4}{*}{\textbf{Cirúrgico (CI)}}
  & CI-1 & Retorno inesperado à sala cirúrgica \\
  & CI-2 & Mudança de procedimento intraoperatório \\
  & CI-3 & Admissão em UTI no pós-operatório \\
  & CI-4 & Readmissão em 30 dias \\
\midrule
\multirow{5}{*}{\textbf{UTI (UT)}}
  & UT-1 & Pneumonia associada à ventilação mecânica \\
  & UT-2 & Reintubação em 24 h \\
  & UT-3 & Úlcera por pressão \\
  & UT-4 & Infecção da corrente sanguínea associada a cateter \\
  & UT-5 & Delirium \\
\midrule
\multirow{4}{*}{\textbf{Perinatal (PE)}}
  & PE-1 & Transfusão sangüínea \\
  & PE-2 & Apgar $< 7$ no 5º minuto \\
  & PE-3 & Admissão em UTI neonatal \\
  & PE-4 & Lesão de parto \\
\midrule
\multirow{3}{*}{\textbf{E.D./Emergência (ED)}}
  & ED-1 & Readmissão em 24 h \\
  & ED-2 & Tempo de permanência $> 6$ h \\
  & ED-3 & Retorno para internação após alta \\
\midrule
\multirow{4}{*}{\textbf{Infecções (IN)}}
  & IN-1 & Infecção do sítio cirúrgico \\
  & IN-2 & Infecção urinária associada a cateter \\
  & IN-3 & \textit{Clostridium difficile} positivo \\
  & IN-4 & Pneumonia adquirida no hospital \\
\midrule
\multirow{4}{*}{\textbf{Geral (GE)}}
  & GE-1 & Transferência para nível de cuidados mais elevado \\
  & GE-2 & Código de ressuscitação cardiopulmonar \\
  & GE-3 & Queda do paciente \\
  & GE-4 & Qualquer procedimento com nota de evento \\
\bottomrule
\end{tabularx}
\fonte{Adaptado de \citeonline{griffin2009ihi}.}
\end{quadro}

% ===========================================================================
\chapter{Casos de Teste Automatizados -- Exemplos de Cobertura}
\label{apend_testes}
% ===========================================================================

Os exemplos a seguir ilustram a estratégia de testes automatizados adotada
no SIGEA-GTT, com cobertura de 100\% de ramos e linhas em ambas as camadas.

\begin{codigo}[htb]
\caption{Exemplo de Teste Unitário -- UserService (JUnit 5 + Mockito).}
\label{cod_teste_user}
\begin{lstlisting}[language=Java]
@ExtendWith(MockitoExtension.class)
class UserServiceTest {

    @Mock
    private UserRepository userRepository;

    @Mock
    private PasswordEncoder passwordEncoder;

    @InjectMocks
    private UserServiceImpl userService;

    @Test
    @DisplayName("Deve rejeitar e-mail fora do dominio @academico.ufs.br")
    void deveLancarExcecaoParaEmailForaDoDominio() {
        CreateUserRequest req = new CreateUserRequest(
            "usuario@gmail.com", UserRole.STUDENT, "123456789");

        assertThrows(DomainEmailException.class,
            () -> userService.createUser(req));

        verify(userRepository, never()).save(any());
    }

    @Test
    @DisplayName("Deve criar usuario STUDENT com matricula obrigatoria")
    void deveCriarUserStudentComMatricula() {
        CreateUserRequest req = new CreateUserRequest(
            "aluno@academico.ufs.br", UserRole.STUDENT, "202100114080");

        when(passwordEncoder.encode(any())).thenReturn("$2a$12$hashed");
        when(userRepository.save(any())).thenAnswer(inv -> inv.getArgument(0));

        UserResponse response = userService.createUser(req);

        assertNotNull(response);
        assertEquals("aluno@academico.ufs.br", response.email());
        assertEquals("202100114080", response.matricula());
    }
}
\end{lstlisting}
\end{codigo}

\begin{codigo}[htb]
\caption{Exemplo de Teste de Integração -- AuthController (Spring Boot Test).}
\label{cod_teste_auth}
\begin{lstlisting}[language=Java]
@SpringBootTest(webEnvironment = RANDOM_PORT)
@AutoConfigureMockMvc
class AuthControllerIT {

    @Autowired MockMvc mockMvc;
    @Autowired ObjectMapper objectMapper;

    @Test
    @DisplayName("POST /auth/login deve retornar JWT para credenciais validas")
    void deveRetornarJwtParaCredenciaisValidas() throws Exception {
        LoginRequest req = new LoginRequest(
            "admin@academico.ufs.br", "Senha@123!");

        mockMvc.perform(post("/api/v1/auth/login")
               .contentType(MediaType.APPLICATION_JSON)
               .content(objectMapper.writeValueAsString(req)))
               .andExpect(status().isOk())
               .andExpect(jsonPath("$.token").isNotEmpty())
               .andExpect(jsonPath("$.expiresAt").isNotEmpty());
    }
}
\end{lstlisting}
\end{codigo}

% ===========================================================================
\chapter{Especificação OpenAPI 3 -- Endpoints Principais}
\label{apend_openapi}
% ===========================================================================

A documentação completa da API REST do SIGEA-GTT é gerada automaticamente
pelo SpringDoc OpenAPI 3 e está disponível em \texttt{/swagger-ui.html}
na aplicação em execução. O \autoref{qua_endpoints} lista os principais
grupos de \textit{endpoints} implementados:

\begin{quadro}[htb]
\centering
\small
\caption{Principais endpoints da API REST do SIGEA-GTT.}
\label{qua_endpoints}
\begin{tabularx}{\textwidth}{l l X}
\toprule
\textbf{Método} & \textbf{Rota} & \textbf{Descrição} \\
\midrule
POST & \texttt{/api/v1/auth/login} & Autenticação e emissão de JWT \\
POST & \texttt{/api/v1/auth/refresh} & Renovação do token de acesso \\
POST & \texttt{/api/v1/auth/change-password} & Troca de senha autenticada \\
\midrule
GET  & \texttt{/api/v1/users} & Listar usuários paginado (10/pág.) \\
POST & \texttt{/api/v1/users} & Cadastrar usuário (gera senha provisória) \\
PUT  & \texttt{/api/v1/users/\{id\}} & Atualizar perfil de usuário \\
DELETE & \texttt{/api/v1/users/\{id\}} & Desativar usuário (soft delete) \\
\midrule
GET  & \texttt{/api/v1/modulos-gtt} & Listar módulos GTT \\
POST & \texttt{/api/v1/modulos-gtt} & Criar módulo GTT \\
GET  & \texttt{/api/v1/gatilhos-gtt} & Listar gatilhos GTT \\
POST & \texttt{/api/v1/gatilhos-gtt} & Criar gatilho GTT \\
\midrule
GET  & \texttt{/api/v1/turmas} & Listar turmas \\
POST & \texttt{/api/v1/turmas} & Criar turma \\
PATCH & \texttt{/api/v1/turmas/\{id\}/encerrar} & Encerrar turma \\
\midrule
GET  & \texttt{/api/v1/atividades} & Listar atividades da turma \\
POST & \texttt{/api/v1/atividades} & Criar atividade com prontuário \\
\midrule
POST & \texttt{/api/v1/submissoes} & Submeter resolução \\
PATCH & \texttt{/api/v1/submissoes/\{id\}/corrigir} & Registrar correção e nota \\
\midrule
GET  & \texttt{/api/v1/dashboard/turma/\{id\}} & Indicadores IHI da turma \\
\bottomrule
\end{tabularx}
\fonte{Elaborado pelo autor (2026).}
\end{quadro}

\end{apendicesenv}
